% ============================================================================
% Sección 4: Implementación del Sistema
% ============================================================================
\section{Implementación del Sistema}

Esta sección describe la arquitectura, diseño e implementación del framework de búsqueda desarrollado. El sistema se compone de tres capas independientes que cooperan para proveer una infraestructura de experimentación reproducible y extensible.

\subsection{Arquitectura general}

El sistema se organiza como un ecosistema de tres repositorios con responsabilidades claramente separadas:

\begin{enumerate}
    \item \textbf{\texttt{search-library}} (Core): Librería publicada en PyPI que implementa los algoritmos de búsqueda como componentes genéricos reutilizables. No posee dependencias externas.

    \item \textbf{\texttt{graph-search}} (Capa experimental -- grafos): Aplicación de benchmark que modela redes de ciudades reales como grafos ponderados y ejecuta comparaciones algorítmicas.

    \item \textbf{\texttt{path-search}} (Capa experimental -- grids): Aplicación de benchmark que evalúa los algoritmos en grids bidimensionales con obstáculos, terrenos variables y múltiples esquemas de movimiento.
\end{enumerate}

La Figura~\ref{fig:system_arch} ilustra la relación entre las capas del sistema.

\begin{figure}[H]
\centering
\begin{tikzpicture}[
    layer/.style={rectangle, draw=blue!60, fill=blue!5, thick, minimum width=6cm, minimum height=1cm, font=\small, text centered},
    dep/.style={-{Stealth[length=3mm]}, thick, dashed},
    node distance=1.2cm
]
    \node[layer] (article) {search-article (Documentación)};
    \node[layer, below=of article, fill=green!5, draw=green!60] (graph) {graph-search (Benchmark Grafos)};
    \node[layer, right=0.3cm of graph, fill=green!5, draw=green!60, minimum width=5.5cm] (path) {path-search (Benchmark Grids)};
    \node[layer, below=1.5cm of graph, fill=orange!5, draw=orange!60, minimum width=12cm] (core) {search-library (Core -- PyPI)};

    \draw[dep] (graph.south) -- ++(0,-0.5) -| (core.north west);
    \draw[dep] (path.south) -- ++(0,-0.5) -| (core.north east);
    \draw[dep] (article.south) -- ++(0,-0.3) -| (graph.north);
    \draw[dep] (article.south) -- ++(0,-0.3) -| (path.north);
\end{tikzpicture}
\caption{Arquitectura del sistema: tres capas con dependencia unidireccional hacia el core.}
\label{fig:system_arch}
\end{figure}

Este diseño garantiza que la librería core permanezca independiente de las capas experimentales, facilitando su reutilización en otros dominios sin arrastrar dependencias de benchmark.

\subsection{Diseño de la librería core (search-library)}

La librería \texttt{search-library} (v1.0.0, publicada en PyPI) constituye el núcleo del sistema. Implementa cinco algoritmos de búsqueda sobre una abstracción genérica parametrizada por tipo, sin dependencias externas de runtime.

\subsubsection{Abstracción SearchProblem}

La interfaz \texttt{SearchProblem[T]} define el contrato que todo problema de búsqueda debe satisfacer. El parámetro genérico \texttt{T} representa el tipo del estado, permitiendo que el mismo algoritmo opere sobre grafos de ciudades (\texttt{T = str}), posiciones en grids (\texttt{T = tuple[int, int]}), o cualquier tipo hashable definido por el usuario:

\begin{lstlisting}[caption={Interfaz abstracta SearchProblem[T]}, label={lst:searchproblem}]
class SearchProblem(ABC, Generic[T]):
    @abstractmethod
    def initial_state(self) -> T: ...
    @abstractmethod
    def is_goal(self, state: T) -> bool: ...
    @abstractmethod
    def successors(self, state: T)
        -> list[tuple[T, float]]: ...
    def heuristic(self, state: T) -> float:
        return 0.0  # Dijkstra por defecto
\end{lstlisting}

El método \texttt{heuristic} provee un valor por defecto de cero, lo que convierte a cualquier algoritmo informado en su variante no informada equivalente (A* con $h=0$ degenera en Dijkstra).

\subsubsection{Jerarquía SearchAlgorithm}

Todos los algoritmos heredan de la clase abstracta \texttt{SearchAlgorithm[T]}, que define una interfaz uniforme con parámetros de control:

\begin{lstlisting}[caption={Base abstracta para algoritmos de búsqueda}, label={lst:searchalg}]
class SearchAlgorithm(ABC, Generic[T]):
    def __init__(self, problem: SearchProblem[T]):
        self._problem = problem

    @abstractmethod
    def search(
        self, *,
        max_iterations: int | None = None,
        strict: bool = False,
        track_explored: bool = False,
    ) -> SearchResult[T]: ...
\end{lstlisting}

Los parámetros \texttt{max\_iterations}, \texttt{strict} y \texttt{track\_explored} permiten controlar el comportamiento sin modificar la lógica algorítmica. En modo \texttt{strict}, las fallas generan excepciones tipadas (\texttt{NoSolutionFoundError}, \texttt{SearchTimeoutError}).

\subsubsection{Algoritmos implementados}

La librería implementa cinco algoritmos con las siguientes características:

\begin{itemize}
    \item \textbf{A* Search (\texttt{AStarSearch}):} Búsqueda informada con cola de prioridad ordenada por $f(n) = g(n) + h(n)$. Emplea \texttt{heapq} con contador para ruptura de empates FIFO y un diccionario \texttt{g\_scores} para detección de caminos mejores.

    \item \textbf{Dijkstra (\texttt{Dijkstra}):} Caso particular de A* con $h(n) = 0$. Implementación independiente con cola de prioridad ordenada exclusivamente por $g(n)$.

    \item \textbf{BFS (\texttt{BFS}):} Búsqueda en amplitud con cola FIFO (\texttt{deque}). Óptima solo en grafos con costos uniformes.

    \item \textbf{DFS (\texttt{DFS}):} Búsqueda en profundidad con pila (LIFO). Sin garantía de optimalidad; retorna el primer camino encontrado.

    \item \textbf{Bidirectional Search (\texttt{BidirectionalSearch}):} BFS simultáneo desde el origen y el destino. Requiere un \texttt{reverse\_problem} explícito que defina sucesores en dirección inversa. Reduce el espacio de búsqueda de $O(b^d)$ a $O(b^{d/2})$.
\end{itemize}

Cada algoritmo expone además una función de conveniencia a nivel de módulo (\texttt{astar\_search}, \texttt{dijkstra\_search}, etc.) que instancia y ejecuta el algoritmo en una sola llamada.

\subsubsection{Estructura SearchResult}

El resultado de toda búsqueda es un \texttt{dataclass} inmutable que encapsula la solución completa junto con métricas de exploración:

\begin{lstlisting}[caption={Contenedor de resultados inmutable}, label={lst:searchresult}]
@dataclass(frozen=True)
class SearchResult(Generic[T]):
    path: list[T]
    total_cost: float
    nodes_explored: int
    success: bool
    explored_states: frozenset[T] | None
    path_length: int  # len(path)
    steps: int        # len(path) - 1
\end{lstlisting}

La inmutabilidad (\texttt{frozen=True}) garantiza que los resultados no sean alterados accidentalmente durante el análisis comparativo.

\subsubsection{Framework de heurísticas}

Las funciones heurísticas se implementan mediante el patrón Strategy a través de la interfaz \texttt{Heuristic[T]}:

\begin{lstlisting}[caption={Patrón Strategy para heurísticas intercambiables}, label={lst:heuristic_iface}]
class Heuristic(ABC, Generic[T]):
    @abstractmethod
    def estimate(self, state: T, goal: T) -> float:
        """Nunca sobreestima el costo real."""

    def __call__(self, state: T, goal: T) -> float:
        return self.estimate(state, goal)
\end{lstlisting}

La librería core provee dos heurísticas para grids:
\begin{itemize}
    \item \textbf{ManhattanHeuristic:} $h(n) = |x_n - x_g| + |y_n - y_g|$. Admisible para movimiento cardinal con costo uniforme.
    \item \textbf{EuclideanHeuristic:} $h(n) = \sqrt{(x_n - x_g)^2 + (y_n - y_g)^2}$. Admisible universalmente, pero menos ajustada que heurísticas especializadas.
\end{itemize}

\subsubsection{Representaciones de dominio}

La librería incluye dos estructuras de datos concretas con sus respectivos adaptadores al \texttt{SearchProblem}:

\begin{itemize}
    \item \textbf{Graph[T]:} Grafo ponderado dirigido o no dirigido con lista de adyacencia. El adaptador \texttt{GraphSearchProblem} traduce la interfaz del grafo a \texttt{SearchProblem[T]}.

    \item \textbf{Grid:} Tablero 2D con soporte para obstáculos, costos variables por celda, y movimiento configurable (4 u 8 direcciones). El adaptador \texttt{GridSearchProblem} acepta una heurística intercambiable.
\end{itemize}

La clase \texttt{Grid} soporta inicialización desde matrices numéricas mediante \texttt{Grid.from\_matrix()}, facilitando la definición declarativa de escenarios experimentales.

\subsection{Capa experimental: graph-search}

El paquete \texttt{graph-city-pathfinding} modela una red vial real del sur de Perú como un grafo ponderado no dirigido, donde los nodos representan ciudades y los pesos de aristas corresponden a distancias aproximadas en kilómetros.

\subsubsection{Modelado del grafo}

La red contiene 10 ciudades (Arequipa, Puno, Juliaca, Cusco, Tacna, Moquegua, Ilo, Nazca, Ayacucho, Abancay) conectadas por 13 aristas con pesos entre 45 km y 560 km. El grafo se diseñó deliberadamente con:

\begin{itemize}
    \item Múltiples rutas entre origen y destino (para evidenciar diferencias algorítmicas).
    \item Una ruta ``trampa'' con menor número de saltos pero mayor costo total (Arequipa $\rightarrow$ Nazca $\rightarrow$ Ayacucho $\rightarrow$ Abancay $\rightarrow$ Cusco: 1600 km), que demuestra la suboptimalidad de BFS en grafos ponderados.
    \item Coordenadas 2D por ciudad para el cálculo de la heurística Euclidiana.
\end{itemize}

\subsubsection{Heurística personalizada}

El módulo implementa \texttt{CityEuclideanHeuristic}, que satisface \texttt{Heuristic[str]} y computa la distancia en línea recta entre ciudades basándose en coordenadas precalculadas. Su admisibilidad se fundamenta en que la distancia Euclidiana es siempre menor o igual a la distancia real por carretera.

\subsubsection{Motor de comparación}

El módulo \texttt{comparison.py} ejecuta los cinco algoritmos sobre el mismo problema, recopilando métricas de costo, nodos explorados, tiempo de ejecución, y optimalidad (verificada contra el baseline de Dijkstra).

\subsection{Capa experimental: path-search}

El paquete \texttt{grid-path-benchmark} evalúa los algoritmos sobre 8 escenarios de grids bidimensionales predefinidos, cubriendo un espectro representativo de configuraciones.

\subsubsection{Sistema de terrenos}

Se define un sistema de terrenos con costos diferenciados que permite modelar entornos realistas:

\begin{itemize}
    \item \texttt{FREE} (costo 1.0): terreno transitable normal.
    \item \texttt{DIFFICULT} (costo 2.0): terreno difícil (arena, barro).
    \item \texttt{WATER} (costo 5.0): terreno acuático (alto costo de travesía).
    \item \texttt{BLOCKED} ($\infty$): obstáculo infranqueable.
\end{itemize}

\subsubsection{Heurísticas extendidas}

Además de las heurísticas provistas por el core, esta capa implementa:

\begin{itemize}
    \item \textbf{ChebyshevHeuristic:} $h(n) = \max(|dx|, |dy|)$. Admisible para grids 8-direccionales con costo uniforme de movimiento diagonal.
    \item \textbf{OctileHeuristic:} $h(n) = \max(dx, dy) + (\sqrt{2} - 1) \cdot \min(dx, dy)$. La heurística más ajustada para grids 8-direccionales donde el movimiento diagonal tiene costo $\sqrt{2}$.
\end{itemize}

La selección automática de heurística se realiza según el escenario: Manhattan para 4 direcciones, Octile para 8 direcciones, y Euclidiana para grids con costos variables.

\subsubsection{Escenarios de benchmark}

Se definen 8 escenarios deterministas con semillas explícitas (\texttt{seed}) que garantizan reproducibilidad total:

\begin{table}[H]
\centering
\caption{Escenarios de benchmark en grids}
\label{tab:scenarios}
\renewcommand{\arraystretch}{1.2}
\begin{tabular}{@{}clcc@{}}
\toprule
\textbf{\#} & \textbf{Escenario} & \textbf{Grid} & \textbf{Movimiento} \\
\midrule
1 & 4-Dir básico (10\% obst.) & $20 \times 20$ & Cardinal \\
2 & 8-Dir básico (10\% obst.) & $20 \times 20$ & Diagonal \\
3 & Obstáculos densos (30\%) & $20 \times 20$ & Cardinal \\
4 & Terreno ponderado & $20 \times 20$ & Diagonal \\
5 & Bug Trap (trampa concava) & $20 \times 20$ & Cardinal \\
6 & Grid grande y disperso & $50 \times 50$ & Diagonal \\
7 & Sin camino posible & $10 \times 10$ & Cardinal \\
8 & Inicio = Meta (trivial) & $10 \times 10$ & Cardinal \\
\bottomrule
\end{tabular}
\end{table}

Los escenarios 5 y 7 representan casos extremos: el Bug Trap evalúa el comportamiento ante heurísticas engañosas (formaciones cóncavas que atrapan a algoritmos greedy), mientras que el escenario sin camino verifica la detección correcta de fallas.

\subsubsection{Generación reproducible de grids}

Los grids se generan mediante funciones deterministas que utilizan \texttt{random.Random(seed)} para la colocación de obstáculos y terrenos. Las posiciones de inicio y meta nunca se bloquean. Para grids ponderados, la distribución de terrenos sigue: 60\% libre, 20\% difícil, 10\% agua, 10\% bloqueado.

\subsection{Integración entre capas}

Ambas capas experimentales dependen exclusivamente de \texttt{search-library} (v$\geq$1.0.0) como dependencia de runtime. La integración se materializa en tres puntos:

\begin{enumerate}
    \item \textbf{Reutilización del \texttt{SearchProblem}:} Ambos benchmarks adaptan sus dominios (grafos de ciudades, grids con terrenos) a la interfaz \texttt{SearchProblem[T]} provista por el core, permitiendo que los mismos algoritmos operen sobre ambos dominios sin modificación.

    \item \textbf{Consistencia de métricas:} El \texttt{SearchResult} estandarizado asegura que las métricas recopiladas (nodos explorados, costo total, éxito) sean directamente comparables entre algoritmos y entre dominios.

    \item \textbf{Extensión sin modificación:} Las heurísticas adicionales (\texttt{CityEuclidean}, \texttt{Chebyshev}, \texttt{Octile}) extienden la interfaz \texttt{Heuristic[T]} del core sin requerir cambios en la librería base, validando la extensibilidad del diseño.
\end{enumerate}

\subsection{Propiedades garantizadas por el sistema}

La arquitectura del sistema garantiza formalmente las siguientes propiedades:

\begin{enumerate}
    \item \textbf{Optimalidad de A*:} Con cualquier heurística que satisfaga la interfaz \texttt{Heuristic[T]} y sea admisible ($h(n) \leq h^*(n)$), \texttt{AStarSearch} encuentra el camino de costo mínimo. Esto se verifica experimentalmente comparando contra el baseline de Dijkstra.

    \item \textbf{Suboptimalidad de BFS en grafos ponderados:} BFS minimiza el número de aristas (hop count), no el costo total. Los experimentos en \texttt{graph-search} demuestran rutas con mayor costo cuando los pesos son heterogéneos.

    \item \textbf{Equivalencia A* = Dijkstra cuando $h(n) = 0$:} El diseño de \texttt{SearchProblem.heuristic()} con retorno por defecto de 0.0 formaliza esta equivalencia. Los resultados experimentales confirman idénticos costos y caminos.

    \item \textbf{Reproducibilidad experimental:} Todas las generaciones de grids y grafos son deterministas. Ejecutar los benchmarks en cualquier ambiente produce los mismos resultados numéricos.

    \item \textbf{Consistencia de métricas:} El uso de un único \texttt{SearchResult} como contenedor de salida para todos los algoritmos garantiza que las comparaciones se realicen sobre métricas homogéneas (nodos explorados, costo, longitud de camino).
\end{enumerate}
